% GENERATED FILE. Edit canonical lessons, then run npm run generate:books.
% canonical-source-hash: fnv1a64:bdbf0b4e92a1b564
% canonical-lessons: SA-C269-media-recall, SA-C269-akasavani, SA-C269-antarjalam, SA-C269-six-channel-practice, SA-C269-media-first-retrieval, SA-C269-media-final-retrieval

\chapter{Six communication words, gently joined}
\label{ch:sa-media-communication}
\hlchaptermodality{269}

\begin{chapteropening}
\textbf{By the end of this chapter:} \emph{I can retrieve six familiar media and communication words from meaning.}

Complete chapter 269 by independently retrieving six media and communication meanings a second time.
\end{chapteropening}

\section[(media-word recall)]{Four communication words are already yours}

\label{lesson:SA-C269-media-recall}



\begin{quote}
Say the Sanskrit for a newspaper or magazine, then a telephone.
\end{quote}

\subsection*{Your turn}
\begin{itemize}
\raggedright
  \item \emph{Recall:} \textbf{\sk{पत्रिका}} — \emph{patrikā} — a newspaper or magazine
  \item \emph{Recall:} \textbf{\sk{दूरभाषः}} — \emph{dūrabhāṣaḥ} — a telephone
  \item \emph{Recall:} \textbf{\sk{दूरदर्शनम्}} — \emph{dūradarśanam} — a television
  \item \emph{Recall:} \textbf{\sk{लेखः}} — \emph{lekhaḥ} — a letter or article
\end{itemize}

\begin{grammarlens}[title={What you've built}]
These four meanings were already in the book. Radio and internet are the two missing members of this small lexical field.
\end{grammarlens}

\begin{tcolorbox}[breakable,colback=teal!4,colframe=teal!35!black,title={Before you move on}]
Say the four familiar words again without looking.
\end{tcolorbox}
\section[\texorpdfstring{\sk{आकाशवाणी} (ākāśavāṇī)}{आकाशवाणी (ākāśavāṇī)}]{\sk{आकाशवाणी} (ākāśavāṇī) — radio}

\label{lesson:SA-C269-akasavani}



\begin{quote}
Recall the four familiar words from the previous lesson.
\end{quote}

\subsection*{You'll want to know: \sk{आकाशवाणी}}
\textbf{\sk{आकाशवाणी}} — \emph{ākāśavāṇī} — “radio.” The \href{https://kosha.sanskrit.today/word/en/radio}{Kosha Sanskrit Today radio entry} explicitly attests this modern use. This tranche teaches \textbf{only the modern radio meaning}. It does not ask you to generalise from the word's older “voice from the sky” sense.

\begin{cousinweb}
The cited entry attests the whole modern term but does not document who coined it or when, so its \textbf{origin remains unknown} here. For memory only, the written form visibly contains \textbf{\sk{आकाश}}, “sky,” and \textbf{\sk{वाणी}}, “voice.” The one meaning to retrieve is “radio.”
\end{cousinweb}

\subsection*{Your turn}
\begin{itemize}
\raggedright
  \item \emph{Say it:} \textbf{\sk{आकाशवाणी}} — radio
  \item \emph{Contrast:} \textbf{\sk{आकाशवाणी}} — radio; \textbf{\sk{दूरभाषः}} — telephone
\end{itemize}

\begin{tcolorbox}[breakable,colback=teal!4,colframe=teal!35!black,title={Before you move on}]
What is the sourced modern radio word in this lesson?
\end{tcolorbox}
\section[\texorpdfstring{\sk{अन्तर्जालम्} (antarjālam)}{अन्तर्जालम् (antarjālam)}]{\sk{अन्तर्जालम्} (antarjālam) — the internet}

\label{lesson:SA-C269-antarjalam}



\begin{quote}
Say the new radio word once from memory.
\end{quote}

\subsection*{You'll want to know: \sk{अन्तर्जालम्}}
\textbf{\sk{अन्तर्जालम्}} — \emph{antarjālam} — “the internet.” The modern computing term is explicitly attested in the \href{https://kosha.sanskrit.today/word/en/Internet}{Kosha Sanskrit Today internet entry} and its \href{https://kosha.sanskrit.today/glossary/computer-internet}{Computer — Internet glossary}. This lesson owns only that sourced modern meaning.

\begin{cousinweb}
The cited entries attest the whole modern term but do not document who coined it or when, so its \textbf{origin remains unknown} here. For memory only, the written form visibly contains \textbf{\sk{अन्तर्}}, “between” or “inter-,” and \textbf{\sk{जाल}}, “net.” That “inter-net” shape is not a lesson in using the internet.
\end{cousinweb}

\subsection*{Your turn}
\begin{itemize}
\raggedright
  \item \emph{Say it:} \textbf{\sk{अन्तर्जालम्}} — the internet
  \item \emph{Contrast:} \textbf{\sk{आकाशवाणी}} — radio; \textbf{\sk{अन्तर्जालम्}} — internet
\end{itemize}

\begin{tcolorbox}[breakable,colback=teal!4,colframe=teal!35!black,title={Before you move on}]
Say the internet word once more without looking.
\end{tcolorbox}
\section[Practice]{Six words, one small field}

\label{lesson:SA-C269-six-channel-practice}



\begin{quote}
Say “radio,” then “internet.”
\end{quote}

\subsection*{Your turn}
Retrieve each word before checking:

\begin{itemize}
\raggedright
  \item press: \textbf{\sk{पत्रिका}} — \emph{patrikā}
  \item telephone: \textbf{\sk{दूरभाषः}} — \emph{dūrabhāṣaḥ}
  \item radio: \textbf{\sk{आकाशवाणी}} — \emph{ākāśavāṇī}
  \item television: \textbf{\sk{दूरदर्शनम्}} — \emph{dūradarśanam}
  \item internet: \textbf{\sk{अन्तर्जालम्}} — \emph{antarjālam}
  \item letter or article: \textbf{\sk{लेखः}} — \emph{lekhaḥ}
\end{itemize}

\begin{grammarlens}[title={What you've built}]
You can now distinguish six meanings in one compact lexical field. This is word retrieval, not evidence that you can use devices, follow news, write correspondence, broadcast, or navigate online.
\end{grammarlens}

\begin{tcolorbox}[breakable,colback=teal!4,colframe=teal!35!black,title={Before you move on}]
Which two words were new in this chapter?
\end{tcolorbox}
\section[(independent media retrieval, first …]{Independent media words: round one}

\label{lesson:SA-C269-media-first-retrieval}



\begin{quote}
Recall which two meanings entered in this chapter.
\end{quote}

\subsection*{Your turn: six English cues}
No Sanskrit model is shown before you answer.

Say the Sanskrit for these six meanings, in order: press, telephone, radio, television, internet, and a letter or article.

\begin{grammarlens}[title={What you've built: boundary}]
You have completed one independent six-word retrieval. That does not yet show device use, news comprehension, correspondence, broadcasting, internet navigation, or digital literacy.
\end{grammarlens}

\begin{tcolorbox}[breakable,colback=teal!4,colframe=teal!35!black,title={Before you move on}]
Retrieve radio and internet once more.
\end{tcolorbox}
\section[(independent media retrieval, second …]{Independent media words: round two}

\label{lesson:SA-C269-media-final-retrieval}



\begin{quote}
Recall the two written-media meanings: press, then a letter or article.
\end{quote}

\subsection*{Your turn: a fresh order}
No Sanskrit model is shown before you answer.

Say the Sanskrit for these meanings in a new order: internet, a letter or article, television, press, radio, and telephone.

\begin{grammarlens}[title={What you've built: boundary}]
You can retrieve six media and communication meanings twice without a visible model. This closes one editorial inventory point. It is not evidence of device use, news comprehension, correspondence, broadcasting, online navigation, digital literacy, human validation, learner attainment, or exam readiness.
\end{grammarlens}

\begin{tcolorbox}[breakable,colback=teal!4,colframe=teal!35!black,title={Before you move on}]
Name the six meanings that belong to this bounded field.
\end{tcolorbox}
